\documentclass[11pt,a4paper]{article}
% Compile with XeLaTeX (fontspec and xeCJK).
\usepackage{fontspec}
\setmainfont{LibertinusSerif}[Extension=.otf,UprightFont=*-Regular,BoldFont=*-Bold,ItalicFont=*-Italic,BoldItalicFont=*-BoldItalic]
\setsansfont{LibertinusSans}[Extension=.otf,UprightFont=*-Regular,BoldFont=*-Bold,ItalicFont=*-Italic]
\usepackage{xeCJK}
\setCJKmainfont{Hiragino Mincho ProN}
\xeCJKDeclareCharClass{Default}{"00B7, "2014, "2018 -> "201F, "2026}
\usepackage[english]{babel}
\usepackage[a4paper,margin=2.5cm,top=2.3cm]{geometry}
\usepackage{graphicx}
\usepackage{xcolor}
\usepackage{titlesec}
\usepackage{fancyhdr}
\usepackage[hidelinks]{hyperref}
\definecolor{accent}{RGB}{31,73,125}
\definecolor{muted}{RGB}{93,104,119}
\hypersetup{colorlinks=true,urlcolor=accent,linkcolor=accent,pdfauthor={Romain Fernex},pdftitle={The origins of Korea and the kingdom of Goguryeo}}
\providecommand{\tightlist}{\setlength{\itemsep}{0pt}\setlength{\parskip}{0pt}}
\providecommand{\ul}[1]{\underline{#1}}
\setcounter{secnumdepth}{0}
\titleformat{\section}{\large\bfseries\color{accent}}{}{0em}{}
\setlength{\parindent}{1.2em}
\setlength{\parskip}{0.35em}
\pagestyle{fancy}
\fancyhf{}
\fancyhead[L]{\small\color{muted}Romain Fernex}
\fancyhead[R]{\small\color{muted}\itshape The origins of Korea and the kingdom of Goguryeo}
\fancyfoot[C]{\small\color{muted}\thepage}
\renewcommand{\headrulewidth}{0pt}
\newcommand{\articlefigure}[2]{%
  \begin{center}
  \includegraphics[width=0.55\linewidth,height=6.5cm,keepaspectratio]{#1}%
  \ifx\relax#2\relax\else\par\smallskip{\small\itshape\color{muted}#2\par}\fi
  \end{center}}

\begin{document}
\thispagestyle{empty}
\noindent{\small\sffamily\color{accent}\MakeUppercase{History of Korea · Part I}}\par\medskip
\noindent{\LARGE\bfseries The origins of Korea and the kingdom of Goguryeo\par}
\smallskip\noindent{\large\itshape Between nationalist narratives and historical realities\par}
\medskip
\noindent{\color{accent}\rule{\linewidth}{0.8pt}}\par\smallskip
\noindent{\small\color{muted}Romain Fernex · 2022 · Article in English}\par\bigskip

``The strongest cause for feelings of nationality ... is identity of political antecedents; the possession of national history, and consequent community of recollections; collective pride and humiliation, pleasure and regret, connected with the same incidents in the past'', John Stuart Mill

In analysing the roots of national feeling in this way, Stuart Mill reminds us of a major, yet sometimes forgotten, role of history: its political role. Indeed, any nation, whether recently formed or not, that seeks to anchor itself in the long term often draws on history, and particularly on its instrumentalisation, as a way of asserting its right, its legitimacy, to exist or develop within a particular geographical or cultural space. The construction of an official history promoted by a government, sometimes far removed from, if not contrary to, established historical realities, thus very often supports claims within the nation itself as well as in its external relations.

In 2004, following applications by China and then North Korea, the archaeological sites of the tombs and capital of the ancient kingdom of Goguryeo{[}1{]}, located on either side of the border between North Korea and China, were inscribed on UNESCO's World Heritage List.

\articlefigure{media/media/image1.jpeg}{Tomb of Gwanggaeto the Great, also known as the General's Tomb, dating from the 4th century BC}

This decision occasioned many controversies, concerning the application that led to it, the process that resulted in the designation, and the reactions of the parties involved after the decision was taken. Thus, the removal of all references to South Korea from Goguryeo's history as presented on the Chinese Ministry of Foreign Affairs website provoked both popular and diplomatic outrage among South Koreans, who felt that their history was being stolen. Beyond traditional diplomatic channels, numerous online forums saw Chinese and Korean participants clash, advancing arguments that were sometimes moderate but often steeped in strong nationalism and historical revisionism. On one side, Chinese internet users maintained that the kingdom of Goguryeo and the populations that preceded it were attached to China and therefore formed part of Chinese cultural unity; on the other, Korean users defended Goguryeo's role as a pillar of Korean identity and as one of the first concrete manifestations of Korea's control over the peninsula.

But what exactly was this kingdom of Goguryeo, and why does it arouse such passions? Throughout this article, we will try to offer some answers by examining, first, Goguryeo's origins, which occupy a particularly important place in this debate, especially on the Chinese side, and, second, the kingdom itself and its importance, not only in the region's historiography but also in claims to a Korean identity stretching back several millennia.

400,000 BC: this is the date of the earliest traces of human populations on the Korean peninsula during the Palaeolithic. These were nomadic peoples, mainly from Manchuria and Northeast Asia, living primarily by hunting and fishing. The Jeulmun period began in the 6th millennium BC, marked by the appearance of comb-pattern pottery found mainly along the Han River, which flows through present-day Seoul.

\articlefigure{media/media/image2.jpeg}{Pottery from the Jeulmun period}

Rice cultivation developed from 4000 BC, aided by a near-tropical summer climate accompanied by heavy monsoons in the same season, allowing some of these peoples to settle and their populations to grow. This development continued until 2000 BC with the transition to the Mumun period, itself marked by a new style of pottery with double-rimmed vessels. This period brought major discoveries conducive to settled life and the emergence of complex, organised civilisations, with the use of bronze in 700 BC and iron in 300 BC. China's growing influence in the region was also characteristic of this period, and it would play an increasing role over the following centuries. It was in this context that the kingdom of Gojoseon emerged in the northern peninsula, itself playing an important role in Korean identity and mythology.

Indeed, for Koreans, Gojoseon's origins lie in the myth of Dangun Waggeon, which dates its creation to 2333 BC, well before the beginning of the Mumun period. According to this myth, Dangun is Hwanung, son of Hwanin, the sky god in Cheondoism and Jeungsanism, two major religious currents in Korea besides Buddhism, which remains the dominant religion and identifies him with Indra, and of Ungnyo, a bear transformed into a woman by Hwanung. He founded the city of Asadal, Gojoseon's capital, on the site of present-day Pyongyang. Although most historians agree that Gojoseon was more akin to a grouping of tribes loosely connected to one another, it nevertheless holds particular importance in Korean culture. This first entity was conquered by Han China, ruled by Emperor Wu, in the 2nd century BC, and then came under its direct control for the first and last time in its history. The Han established four commanderies as administrative divisions to govern Gojoseon's territory.

\articlefigure{media/media/image3.jpeg}{The four commanderies during the period of Chinese rule}

The commanderies of Lelang and Xuantu were the most important, with Lelang producing textiles and ceramics, skills that Koreans would preserve and develop throughout their history. This was therefore an important period of Sinicisation in the region, giving rise to numerous connections and exchanges between China and the populations of the Korean peninsula. It also provides substantial support for Chinese claims to include Korea within the sphere of Chinese civilisation. However, while China undeniably played a major role in this region at the time, this overlooks the many populations, clans and tribes that remained there and with which China maintained relations that were sometimes hostile, sometimes friendly. China indeed pursued a policy of appeasement, also called Heqin, consisting of showering certain clans with gifts to win their favour.

In the south were three Han tribes: the Mahan, the most numerous and based on agriculture; the Chinhan, dependent particularly on silk production; and the Pyonhan. These would give rise to several kingdoms clearly distinct from China a few centuries later, which we will have the opportunity to revisit in another article. Alongside these three tribes were the Wa, the first population associated with Japan, who nevertheless retained a significant presence in the southeastern Korean peninsula.

Finally, in the north, on Gojoseon's territory, there were a total of nine tribes. Among them were Buyeo, a close ally of China, and Goguryeo. The latter appeared in 37 BC according to the \emph{Samguk Sagi} {[}2{]} (History of the Three Kingdoms), founded by Dongmyeong in southern Manchuria.

\articlefigure{media/media/image4.jpeg}{King Dongmyeong of Goguryeo (58 BC -- 19 BC)}

A particularly warlike tribal grouping because its territory lacked sufficient resources to sustain its growing population, it rapidly conquered other territories and organised itself into four major administrative divisions. It even maintained diplomatic relations with China from 75 to 12 BC, and its ruler proclaimed himself king (wang in Korean) to emphasise his desire to distinguish himself from China. Indeed, although China sought the favour of certain clans, it also subjected many of the region's native populations to severe exploitation, provoking frequent revolts. Goguryeo's meteoric rise therefore worried China all the more because it was gaining influence among other local tribes opposed to China.

The breakup of the Han Empire in AD 220 gave many powers, including Goguryeo, the opportunity to free themselves from Chinese influence and develop an initial sense of Korean identity, or rather identities, given the still relatively heterogeneous character of these different tribes, despite their many similarities. This rise was short-lived with the ascent of China's Wei dynasty in northern China, which, fearing Goguryeo and wishing to pacify the region, attacked it and took control of its capital in 244. However, Goguryeo's disappearance was only temporary, a prelude to a rebirth barely a century later that would establish it as one of the peninsula's most powerful kingdoms for several centuries.

It is to this period of intense rivalries between emerging powers on the peninsula, whose three largest, Goguryeo, Baekje and Goguryeo, would confront one another in armed conflicts and struggles for influence that profoundly marked the region's history, that we will return in the next article in this series. See you soon!

\textbf{Bibliography:}

\begin{itemize}
\item
  Michael J.Seth, \emph{A history of Korea -- From Antiquity to the present}, 2010, Rowman \& Littlefield Publishers
\item
  Kyung Moon Hwang, \emph{A History of Korea}, 3rd edition, 2021, Palgrave
\item
  Kim Taesik, \emph{Koguryeo and Gaya - contacts and consequences}, Hungik universitt, journal of Northeast Asian History, Vol. 4 n°1 (p41-82), 2007
\item
  Jung Woon Yong, \emph{Trends in Koguryo's relationship with Paekche and Silla during the 4th -- 7th centuries}, International Journal of Korean history, Vol. 8, 2006
\item
  Mark Byington, \emph{The War of Words Between South Korea and China Over An Ancient Kingdom : Why Both sides are misguided}, history news network, Columbian college of arts \& science, The George Washington university, 2004
\item
  Thomas Chase, \emph{Nationalism and the Net} \emph{: online discussion of Goguryeo history in China and South Korea}, Monash university, China Information, 2011
\item
  John Stuart Mill, Considerations on representative government, Ch. XVI, Of Nationality, as connected with representative government, 1861
\end{itemize}

\textbf{{[}1{]}} For consistency, I use throughout this article the Revised Romanisation of Korean, which has been Korea's official romanisation system since the 2000s and is also used by Kyung Moon Hwang. However, a substantial proportion, if not the majority, of the articles I used in writing this piece employ different spellings; for instance, Michael J. Seth uses the McCune-Reischauer system, which was authoritative in English-language academia when he wrote his book. Thus, Goguryeo may appear as Koguryo ̆ and, more generally, the names of places, kingdoms, key figures and so on are spelt differently from those in this article.

\textbf{{[}2{]}} The \emph{Samguk Sagi}, or History of the Three Kingdoms in Korean, is one of the major works, alongside the \emph{Samguk Yusa} (Memorabilia of the Three Kingdoms), published in 1279, dealing with the region's history in this period, reaching back several centuries before our era. Written in 1145 by the Korean historian and general Kim Busik, it is indeed one of the main written sources on the Korean peninsula available to historians of this period\ldots{} Yet certain passages are particularly romanticised and strongly influenced by the period in which the text was written, for reasons we will probably revisit in another series of articles, a period very different from the one it describes. Moreover, the foundational text of Japanese history, the \emph{Nihon Shoki}, and several major Chinese texts, including the \emph{Shiji} by the Chinese historian Sima Qian, himself a contemporary of this period, published in 100 BC, are also valuable sources for documenting it.

\end{document}
